\uuid{K4Wy}
\titre{Dérivée en un point}
\auteur{bodin ; hanani ; mzari}
\organisation{exo7}
\licence{CC BY-NC-SA 4.0}
\datecreate{2020-01-20}
\niveau{L1}
\module{Analyse}
\chapitre{Dérivabilité des fonctions réelles}
\sousChapitre{Autre}
\difficulte{2}
\theme{Dérivées, dérivée en un point}
\exo7qcmid{lille-1.derivees.014}

\contenu{
\begin{question}
\qclassification{124.00}
\qtags{motcle=dérivée en un point}
Soit $a,b\in \Rr$ et $f(x)=\left\{\begin{array}{cl}\displaystyle \mathrm{e}^{x^2+x}&\mbox{si }x\leq 0\\ \\ a \arctan x+b &\mbox{si }x>0.\end{array}\right.$

Quelles valeurs faut-il donner à $a$ et $b$ pour que $f$ soit dérivable sur $\Rr$ ?
\begin{answers}
    \bad{$a=1$ et $b=0$}
    \bad{$a=0$ et $b=1$}
    \bad{$a=0$ et $b=0$}
    \good{$a=1$ et $b=1$}
\end{answers}
\begin{explanations}
La fonction $f$ est continue sur $\Rr^*$ et pour qu'elle soit continue en $0$, il faut que
$$\lim _{x\to 0^-}f(x)=\lim _{x\to 0^+}f(x)\Rightarrow 1=b.$$
De plus, $f$ est dérivable sur $\Rr^*$ avec
$$f'(x)=\left\{\begin{array}{cl}\displaystyle (2x+1)\mathrm{e}^{x^2+x}&\mbox{si }x<0\\ \\ \displaystyle \frac{a}{1+x^2} &\mbox{si }x>0\end{array}\right.$$
et $f'_g(0)=1$ et $f'_d(0)=a$. Donc, pour que $f$ soit dérivable en $0$, on doit avoir $f'_g(0)=f'_d(0)$. D'où $a=1$.
\end{explanations}
\end{question}
}
